% ---
% Informações de dados para CAPA e FOLHA DE ROSTO
% ---
\titulo{Produto 2 - Relatório contendo proposta de solução para gerar relatório agregado que consolide e organize as contribuições recebidas através ferramentas e aplicações de consulta pública}
\autor{Diego Rabatone Oliveira}
\local{Brasil}
\data{Abril/2015, v1}
\newcommand{\numeroContrato}{2015/000042-00}
\newcommand{\valorProduto}{R\$6.200,00 (seis mil e duzentos reais)}
\newcommand{\valorContrato}{R\$40.000,00 (quarenta mil reais}
\newcommand{\dataEntrega}{08/03/2015}
\newcommand{\labelSupervisor}{Supervisora}
\newcommand{\supervisor}{Sabrina Durigon Marques}
\instituicao{%
  Secretaria de Assuntos Legislativos do 
  \par
  Ministério da Justiça (SAL/MJ)}
\orientador[Supervisora:]{Sabrina Durigon Marques}

\tipotrabalho{Relatório técnico}
% O preambulo deve conter o tipo do trabalho, o objetivo, 
% o nome da instituição e a área de concentração 
\preambulo{Relatório referente ao Projeto BRA/07/004 – Democratização de Informações no Processo de Elaboração Normativa}
% ---

% ---
% Configurações de aparência do PDF final

% alterando o aspecto da cor azul
\definecolor{blue}{RGB}{41,5,195}

% informações do PDF
\makeatletter
\hypersetup{
     	%pagebackref=true,
		pdftitle={\@title}, 
		pdfauthor={\@author},
    	pdfsubject={\imprimirpreambulo},
	    pdfcreator={LaTeX with abnTeX2},
		pdfkeywords={wordpress}{marco civil da internet}{proteção de dados pessoais}{métodos ágeis}, 
		colorlinks=true,       		% false: boxed links; true: colored links
    	linkcolor=blue,          	% color of internal links
    	citecolor=blue,        		% color of links to bibliography
    	filecolor=magenta,      		% color of file links
		urlcolor=blue,
		bookmarksdepth=4
}
\makeatother
% --- 

% --- 
% Espaçamentos entre linhas e parágrafos 
% --- 

% O tamanho do parágrafo é dado por:
\setlength{\parindent}{1.2cm}

% Controle do espaçamento entre um parágrafo e outro:
\setlength{\parskip}{0.2cm}  % tente também \onelineskip

%%criar um novo estilo de cabeçalhos e rodapés
\makepagestyle{estilosalmj}
  %%cabeçalhos
  \makeevenhead{estilosalmj} %%pagina par
     {\includegraphics[scale=0.2]{./imagens/topo-esquerda.png}}
     {\thepage}
     {\includegraphics[scale=0.4]{./imagens/pnud.png}}
  \makeoddhead{estilosalmj} %%pagina ímpar ou com oneside
     {\includegraphics[scale=0.2]{./imagens/topo-esquerda.png}}
     {}
     {\includegraphics[scale=0.4]{./imagens/pnud.png} %
     \ifnum\value{page}>4 \ABNTEXfontereduzida Página \thepage \else \fi}
  \makeheadrule{estilosalmj}{\textwidth}{\normalrulethickness} %linha
  % rodapé
  \makeevenfoot{estilosalmj}
     {} %%pagina par
     {}
     {}
  \makeoddfoot{estilosalmj} %%pagina ímpar ou com oneside
     {}
     {}
     {}